\documentclass[11pt]{article}
\usepackage[brazilian]{babel}
\newcommand{\imgdir}{images/pt}
% Shared layout for the Illuminati manual. Each manual-<lang>.tex loads this
% after \documentclass and sets \manuallang (babel name) and \imgdir.
% Build with XeLaTeX (see build.ps1); the fonts are the app's own.

\usepackage{fontspec}
\setmainfont{Literata}[
  Path=../../assets/fonts/literata/, Extension=.ttf,
  UprightFont=*-Regular, BoldFont=*-Bold, ItalicFont=*-Italic, BoldItalicFont=*-BoldItalic]
\setsansfont{AtkinsonHyperlegible}[
  Path=../../assets/fonts/atkinson/, Extension=.ttf,
  UprightFont=*-Regular, BoldFont=*-Bold, ItalicFont=*-Italic, BoldItalicFont=*-BoldItalic]
\newfontfamily\dyslexicfont{OpenDyslexic}[
  Path=../../assets/fonts/opendyslexic/, Extension=.otf,
  UprightFont=*-Regular, BoldFont=*-Bold, ItalicFont=*-Italic, BoldItalicFont=*-Bold-Italic]
\newfontfamily\garamondfont{EBGaramond}[
  Path=../../assets/fonts/ebgaramond/, Extension=.ttf,
  UprightFont=*-Regular, BoldFont=*-Bold, ItalicFont=*-Italic, BoldItalicFont=*-BoldItalic]

\usepackage[a4paper, margin=2.4cm, headheight=14pt]{geometry}
\usepackage{graphicx}
\usepackage[table]{xcolor}
\usepackage{booktabs, tabularx, array}
\usepackage{enumitem}
\usepackage{fancyhdr}
\usepackage{titlesec}
\usepackage[most]{tcolorbox}
\usepackage[hidelinks]{hyperref}

% Colours from the app's logo.
\definecolor{violet}{HTML}{34206A}
\definecolor{violetlight}{HTML}{EDE8F7}
\definecolor{gold}{HTML}{C9A227}
\definecolor{goldlight}{HTML}{FBF4DF}

\setlength{\parskip}{0.55em}
\setlength{\parindent}{0pt}
\setlist{itemsep=0.25em, topsep=0.3em}
\linespread{1.08}

\titleformat{\section}{\sffamily\Large\bfseries\color{violet}}{\thesection}{0.8em}{}
\titleformat{\subsection}{\sffamily\large\bfseries\color{violet}}{\thesubsection}{0.7em}{}
\titlespacing*{\section}{0pt}{1.6em}{0.6em}
\titlespacing*{\subsection}{0pt}{1.1em}{0.35em}
% Don't leave a heading alone at the bottom of a page.
\usepackage{needspace, etoolbox}
\pretocmd{\section}{\needspace{8\baselineskip}}{}{}
\pretocmd{\subsection}{\needspace{5\baselineskip}}{}{}

\pagestyle{fancy}
\fancyhf{}
\fancyhead[L]{\sffamily\small\color{violet} Illuminati}
\fancyhead[R]{\sffamily\small\color{gray}\nouppercase{\rightmark}}
\fancyfoot[C]{\sffamily\small\thepage}
\renewcommand{\headrulewidth}{0.3pt}
% \rightmark is the first section mark on the page.
\renewcommand{\sectionmark}[1]{\markboth{#1}{#1}}
\renewcommand{\subsectionmark}[1]{}

% Keep catalogue names whole in every language.
\babelhyphenation{Wiki-source Guten-berg}

% A key or key combination, e.g. \kbd{Ctrl+F}.
\newcommand{\kbd}[1]{\tikz[baseline=(k.base)]\node[draw=gray!60, fill=gray!8,
  rounded corners=2pt, inner sep=2pt, font=\sffamily\small](k){#1};}
\usepackage{tikz}

% Arrow keys (the sans font has no arrow characters).
\newcommand{\kright}{\kbd{\ensuremath{\rightarrow}}}
\newcommand{\kleft}{\kbd{\ensuremath{\leftarrow}}}
\newcommand{\kup}{\kbd{\ensuremath{\uparrow}}}
\newcommand{\kdown}{\kbd{\ensuremath{\downarrow}}}

% A button, menu item or label as it appears in the app.
\newcommand{\ui}[1]{\textsf{\textbf{#1}}}

% Screenshot with caption; skipped if the image file isn't there.
\newcommand{\shot}[2]{%
  \IfFileExists{\imgdir/#1.jpg}{%
    \begin{figure}[htbp]\centering
      \fcolorbox{gray!40}{white}{\includegraphics[width=0.96\textwidth]{\imgdir/#1.jpg}}
      \caption{#2}
    \end{figure}}{}}

\newtcolorbox{tip}[1]{enhanced, colback=goldlight, colframe=gold,
  boxrule=0.6pt, arc=3pt, left=8pt, right=8pt, top=5pt, bottom=5pt,
  fonttitle=\sffamily\bfseries, coltitle=black, colbacktitle=gold!35,
  title={#1}}
\newtcolorbox{note}[1]{enhanced, colback=violetlight, colframe=violet,
  boxrule=0.6pt, arc=3pt, left=8pt, right=8pt, top=5pt, bottom=5pt,
  fonttitle=\sffamily\bfseries, coltitle=white, colbacktitle=violet,
  title={#1}}

% Two-column tables that fill the text width.
\newenvironment{keytable}[2]{%
  \renewcommand{\arraystretch}{1.3}%
  \rowcolors{2}{white}{violetlight}%
  \tabularx{\textwidth}{@{}>{\raggedright\arraybackslash}p{0.33\textwidth}X@{}}
  \toprule \textsf{\textbf{#1}} & \textsf{\textbf{#2}} \\ \midrule}{%
  \bottomrule\endtabularx}

% Title page: \maketitlepage{subtitle}{version line}{blurb}
\newcommand{\maketitlepage}[3]{%
  \begin{titlepage}
    \centering
    \vspace*{3cm}
    \includegraphics[width=4.2cm]{../../assets/icon/app_icon.png}\par\vspace{1.2cm}
    {\sffamily\fontsize{40}{46}\selectfont\bfseries\color{violet} Illuminati\par}
    \vspace{0.5cm}
    {\sffamily\LARGE #1\par}
    \vspace{1.5cm}
    {\large\itshape #3\par}
    \vfill
    {\sffamily\color{gray} #2\par}
  \end{titlepage}}

\hypersetup{pdftitle={Illuminati — Manual do usuário}, pdfauthor={Illuminati}, pdflang={pt}}

\begin{document}

\maketitlepage{Manual do usuário}{Versão 1.0 · Windows, macOS, Linux e Android}%
  {Encontre, baixe e leia livros que nunca tiveram direitos autorais\\ou cujos direitos já expiraram.}

{\setcounter{tocdepth}{1}\tableofcontents}
\newpage

% ---------------------------------------------------------------------------
\section{Boas-vindas}

Illuminati é um aplicativo de leitura para Windows, macOS, Linux e Android. Funciona como os
aplicativos de e-books que você talvez já conheça, com uma diferença: todos os livros que ele
oferece são gratuitos e legais de baixar, porque seus direitos autorais expiraram ou nunca
existiram.

\begin{itemize}
  \item Mais de 75.000 livros do \textbf{Project Gutenberg}, além de \textbf{Standard Ebooks},
        \textbf{Wikisource} e \textbf{Internet Archive}.
  \item Um leitor confortável com destaques, notas, marcadores, busca, dicionário e tradução.
  \item Sua própria biblioteca, com estantes e listas de leitura.
  \item Ajustes de tela e acessibilidade, incluindo uma fonte pensada para a dislexia.
  \item Sem conta, sem assinatura e sem anúncios. Tudo fica no seu dispositivo.
\end{itemize}

A interface está disponível em português, inglês, espanhol, francês, alemão e italiano.

% ---------------------------------------------------------------------------
\section{Primeiros passos}

\subsection{Instalação}
\begin{itemize}
  \item \textbf{Windows:} copie a pasta \texttt{Illuminati} inteira para o computador e execute
        \texttt{illuminati.exe}. Deixe os outros arquivos ao lado dele; o programa precisa deles.
  \item \textbf{Android:} instale o arquivo \texttt{.apk}. Talvez seja preciso permitir a
        instalação de aplicativos desta origem nas configurações do celular.
  \item \textbf{macOS e Linux:} instale a versão preparada para o seu sistema.
\end{itemize}

\subsection{A tela inicial}
No computador, o aplicativo abre maximizado. A tela inicial tem duas partes:

\begin{enumerate}
  \item \ui{Sua biblioteca} (em cima): os livros que você baixou, começando pelo último lido, cada
        um com uma barra de progresso. O botão \ui{Biblioteca e estantes} abre toda a sua
        biblioteca.
  \item \ui{Descubra livros grátis} (embaixo): uma caixa de busca, o idioma do catálogo, as
        categorias, uma fileira para cada catálogo adicional e, por fim, o catálogo do Project
        Gutenberg.
\end{enumerate}

\shot{home}{A tela inicial.}

Os dois botões no canto superior direito abrem \ui{Tela e acessibilidade}
(seção~\ref{sec:display}) e \ui{Sobre}, que mostra de onde vêm os livros e as licenças das fontes.

% ---------------------------------------------------------------------------
\section{Encontrar livros}

\subsection{Busca}
Digite em \ui{Pesquisar por título ou autor}. A busca:
\begin{itemize}
  \item encontra qualquer parte de um título ou do nome de um autor: \emph{miserá} encontra
        \emph{Os Miseráveis};
  \item ignora acentos e maiúsculas: \emph{miseraveis} e \emph{MISERÁVEIS} dão os mesmos
        resultados;
  \item pesquisa todos os catálogos ao mesmo tempo; cada um mostra os resultados na própria
        fileira.
\end{itemize}
Limpe a caixa com o botão \ui{×} para voltar aos livros populares.

\subsection{Categorias e idioma}
As etiquetas abaixo da caixa de busca (\ui{Ficção}, \ui{Mistério}, \ui{Ficção científica},
\ui{Poesia}…) limitam todos os catálogos a um tema. O menu de idioma à direita mostra só os livros
de um idioma, ou \ui{Todos os idiomas}. Os catálogos sem nada para o idioma ou a categoria
escolhida ficam ocultos.

\subsection{Os catálogos}
\begin{keytable}{Catálogo}{O que oferece}
  Project Gutenberg & A maior biblioteca gratuita, com mais de 75.000 livros ordenados por
                      popularidade. Muitos têm edição ilustrada. \\
  Standard Ebooks   & Clássicos revisados com cuidado e bem diagramados, em inglês. \\
  Wikisource        & Textos transcritos por voluntários, em muitos idiomas. O e-book é
                      preparado no momento do download; obras muito longas (por exemplo,
                      romances em vários volumes) podem não estar disponíveis. \\
  Internet Archive  & Digitalizações de edições impressas antigas, publicadas há 95 anos ou mais.
                      O texto foi extraído automaticamente das digitalizações e pode conter
                      erros. \\
\end{keytable}

\subsection{A página de um livro}
Toque em uma capa para abrir a página do livro. Ela mostra o resumo, os temas, o idioma, o
catálogo e uma nota sobre direitos autorais. Dali você pode:
\begin{itemize}
  \item \ui{Baixar grátis}: o livro é salvo no seu dispositivo, pronto para ler sem internet. Cada
        capa do catálogo também tem um pequeno botão de download.
  \item \ui{Ler agora} ou \ui{Continuar lendo}, depois de baixado.
  \item \ui{Quero ler}, \ui{Lido} e \ui{Adicionar a uma estante}, mesmo antes de baixar
        (seção~\ref{sec:library}).
  \item \ui{Obter a edição ilustrada}, para livros do Gutenberg baixados sem imagens.
\end{itemize}

\shot{details}{A página de um livro.}

% ---------------------------------------------------------------------------
\section{Sua biblioteca}\label{sec:library}

\ui{Biblioteca e estantes}, na tela inicial, abre todos os livros que você guardou: os baixados e
os que você salvou para depois sem baixar (com a capa esmaecida e uma nuvenzinha).

\shot{library}{A biblioteca, com suas visualizações e estantes.}

\subsection{Visualizações, filtro e ordem}
As etiquetas no topo escolhem uma visualização, cada uma com o número de livros: \ui{Todos},
\ui{Baixados}, \ui{Quero ler}, \ui{Lendo}, \ui{Lido} e depois cada uma das suas estantes. A caixa
acima filtra por título ou autor. O botão de ordenar (canto superior direito) ordena os livros por
\ui{Lidos recentemente}, \ui{Título}, \ui{Autor} (pelo sobrenome), \ui{Progresso} ou
\ui{Adicionados recentemente}. O aplicativo lembra a sua escolha.

\subsection{Listas de leitura}
\begin{itemize}
  \item \ui{Quero ler} guarda um livro para depois, sem baixá-lo. Quando você abre o livro, ele sai
        dessa lista.
  \item \ui{Lido} é marcado sozinho quando você chega ao fim de um livro (dá para desfazer), ou à
        mão.
\end{itemize}

\subsection{Estantes}
Estantes são suas próprias coleções, como \emph{Favoritos} ou \emph{Filosofia}. Um livro pode
estar em várias estantes.
\begin{itemize}
  \item Crie uma com \ui{Nova estante} na biblioteca, ou em \ui{Adicionar a uma estante} em
        qualquer livro.
  \item Clique com o botão direito (ou toque e segure) em uma estante para renomeá-la ou
        excluí-la. Ao excluir uma estante, os livros continuam na sua biblioteca.
\end{itemize}

\subsection{O menu de um livro}
Clique com o botão direito em um livro (ou toque e segure no celular) para ver
\ui{Detalhes do livro}, \ui{Adicionar a uma estante}, \ui{Quero ler}, \ui{Marcar como lido} e duas
formas de removê-lo:
\begin{itemize}
  \item \ui{Remover o download} libera espaço, mas mantém o livro na biblioteca com notas,
        estantes e posição de leitura.
  \item \ui{Remover da biblioteca} apaga tudo, inclusive notas e marcadores.
\end{itemize}

% ---------------------------------------------------------------------------
\section{Ler}

Toque em um livro baixado para abri-lo. Ele abre onde você parou.

\shot{reader}{O leitor com sua barra de ferramentas.}

\subsection{Virar páginas}
\begin{keytable}{Ação}{Como}
  Próxima página       & Toque na borda direita · \kright{} · \kbd{Page Down} · \kbd{Espaço} \\
  Página anterior      & Toque na borda esquerda · \kleft{} · \kbd{Page Up} \\
  Rolar um pouco       & \kup{} \kdown{} ou a roda do mouse \\
  Mostrar a barra      & Toque no centro da página \\
\end{keytable}

\subsection{A barra de ferramentas}
Toque no centro da página para mostrá-la. Em cima: voltar à biblioteca, busca
(seção~\ref{sec:search}), marcador, sumário e ajustes de tela. Embaixo: capítulo anterior e
seguinte, um controle deslizante para pular entre capítulos e o seu progresso.

\subsection{Progresso e tempo restante}
A linha de baixo mostra quantos minutos faltam no capítulo e quanto do livro você já leu. O tempo
é calculado com a sua própria velocidade de leitura, que o aplicativo aprende enquanto você lê.

\subsection{Sumário e imagens}
\ui{Sumário} mostra os capítulos reais do livro, mesmo quando um único arquivo contém vários. As
imagens aparecem no texto; toque nelas para vê-las em tela cheia e ampliá-las.

% ---------------------------------------------------------------------------
\section{Destaques, notas e marcadores}

\subsection{Selecionar texto}
No celular, toque e segure uma palavra e arraste as alças. No computador, arraste o mouse sobre o
texto ou clique duas vezes em uma palavra e depois clique com o \textbf{botão direito} na seleção.
Aparece um menu com:
\begin{itemize}
  \item \ui{Destacar}: marca o texto com a sua cor de destaque.
  \item \ui{Nota}: um destaque com uma nota sua, em uma de quatro cores.
  \item \ui{Definir} e \ui{Traduzir} (seção~\ref{sec:lookup}).
  \item \ui{Copiar}.
\end{itemize}
Os destaques se ajustam a palavras inteiras. Os que têm nota aparecem sublinhados. Selecione de
novo um trecho destacado para editar a nota ou removê-lo.

\subsection{Marcadores}
Use o botão de marcador da barra de ferramentas, ou \kbd{Ctrl+B}. Uma fita vermelha indica que a
página tem um marcador.

\subsection{O painel lateral}
O botão de sumário abre um painel com três abas: \ui{Sumário}, \ui{Marcadores} e \ui{Notas}.
Toque em qualquer item para ir até ele. \ui{Copiar tudo em Markdown} copia todos os seus destaques
e notas, prontos para colar em outro programa.

% ---------------------------------------------------------------------------
\section{Dicionário e tradução}\label{sec:lookup}

Selecione uma palavra (até três) e escolha \ui{Definir} para ver seus significados. Os do idioma
do livro aparecem primeiro. Eles vêm do Wiktionary e estão escritos em inglês.

Selecione qualquer trecho e escolha \ui{Traduzir}. Escolha o idioma de destino na lista; o
aplicativo lembra dele.

\begin{note}{É preciso internet}
Definições e traduções precisam de conexão com a internet. O texto selecionado é enviado ao
Wiktionary ou ao MyMemory; nada mais é enviado.
\end{note}

% ---------------------------------------------------------------------------
\section{Buscar dentro de um livro}\label{sec:search}

Pressione \kbd{Ctrl+F} ou o botão de busca da barra de ferramentas e digite pelo menos duas
letras. A busca ignora maiúsculas e acentos. Escolha um resultado para ir até ele. Uma barra no
topo mostra \emph{3 de 27}: use as setas dela, ou \kbd{F3}, para ir ao próximo resultado.
\kbd{Esc} a fecha.

% ---------------------------------------------------------------------------
\section{Tela e acessibilidade}\label{sec:display}

Abra esses ajustes com o botão de configurações da tela inicial, ou com o botão de texto do
leitor. Uma prévia no topo mostra cada mudança na hora.

\shot{settings}{Tela e acessibilidade.}

\subsection{Idioma}
\ui{Idioma do app} muda a interface para português, inglês, espanhol, francês, alemão ou italiano
na hora. \ui{Igual ao dispositivo} segue o idioma do seu dispositivo.

\subsection{Texto}
\begin{itemize}
  \item Tema: \ui{Claro}, \ui{Sépia} ou \ui{Escuro}.
  \item Fonte: quatro fontes incluídas no aplicativo, para que os livros tenham a mesma aparência
        em todos os dispositivos:
  \begin{itemize}
    \item \textbf{Literata}, criada para leitura na tela (a padrão);
    \item {\garamondfont EB Garamond}, um tipo clássico de livro;
    \item {\sffamily Atkinson Hyperlegible}, feita para que letras parecidas se distingam com
          facilidade;
    \item {\dyslexicfont OpenDyslexic}, pensada para leitores com dislexia.
  \end{itemize}
  \item \ui{Tamanho do texto}, \ui{Espaçamento entre linhas} e \ui{Espaçamento entre letras}.
  \item \ui{Justificar o texto}: desative para que os espaços entre palavras fiquem iguais, o que
        muitos leitores com dislexia acham mais fácil.
\end{itemize}

\subsection{Página}
\ui{Margens} e \ui{Largura da linha}. Linhas mais curtas são mais fáceis de acompanhar.

\subsection{Tela}
\ui{Brilho} escurece todo o aplicativo. \ui{Temperatura} reduz a luz azul, como um modo noturno.
Os dois se somam aos ajustes do seu dispositivo.

\ui{Restaurar padrões} volta todos esses ajustes ao estado inicial.

% ---------------------------------------------------------------------------
\section{Atalhos de teclado}

\begin{keytable}{Teclas}{Ação}
  \kright{} \kbd{Page Down} \kbd{Espaço} & Próxima página \\
  \kleft{} \kbd{Page Up}                 & Página anterior \\
  \kup{} \kdown{}                        & Rolar um pouco \\
  \kbd{Ctrl+F}                           & Buscar no livro \\
  \kbd{F3}                               & Próximo resultado \\
  \kbd{Ctrl+B}                           & Pôr ou tirar um marcador na página \\
  \kbd{Esc}                              & Fechar a busca ou a barra, ou voltar à biblioteca \\
  \kbd{Shift} + setas                    & Ampliar uma seleção de texto \\
\end{keytable}

% ---------------------------------------------------------------------------
\section{Privacidade e seus dados}

O Illuminati não tem contas nem rastreamento. Sua biblioteca, suas notas e seus ajustes ficam
salvos apenas no seu dispositivo (no Windows, em \texttt{\%APPDATA\%\textbackslash Illuminati}).

O aplicativo só se conecta à internet para:
\begin{itemize}
  \item navegar e pesquisar nos catálogos e baixar livros e capas;
  \item definir ou traduzir o texto que você selecionou, que é enviado ao Wiktionary ou ao
        MyMemory.
\end{itemize}

% ---------------------------------------------------------------------------
\section{Sobre direitos autorais}

Todos os livros estão em domínio público nos Estados Unidos. Alguns textos do Wikisource têm, em
vez disso, licenças livres (como CC BY-SA), que permitem lê-los e compartilhá-los livremente. Do
Internet Archive, só aparecem livros publicados há 95 anos ou mais.

\begin{tip}{Confira as leis do seu país}
As leis de direitos autorais variam de país para país. Um livro livre nos Estados Unidos pode
ainda estar protegido onde você mora, então confira as leis do seu país antes de baixar.
\end{tip}

% ---------------------------------------------------------------------------
\section{Solução de problemas}

\begin{keytable}{Mensagem}{O que fazer}
  Não foi possível acessar um catálogo & O serviço pode estar sobrecarregado ou sua conexão
                             caiu. Use \ui{Tentar novamente} mais tarde; os outros catálogos
                             continuam funcionando. \\
  Este título não tem edição em e-book & Provavelmente é um audiolivro. Procure outra
                             edição. \\
  O Wikisource não conseguiu preparar este livro & A obra é longa demais para o serviço do
                             Wikisource. Tente um dos volumes ou outro catálogo. \\
  Limite diário de traduções gratuitas atingido & O MyMemory permite um número limitado de
                             traduções gratuitas por dia. Tente amanhã. \\
  O download demorou demais & Verifique sua conexão e tente de novo. \\
\end{keytable}

% ---------------------------------------------------------------------------
\section{Créditos}

Livros: Project Gutenberg, Standard Ebooks, Wikisource (via WS Export) e Internet Archive.
Definições: Wiktionary (CC BY-SA). Traduções: MyMemory.

Fontes com a licença SIL Open Font License: Literata, EB Garamond, Atkinson Hyperlegible (Braille
Institute) e OpenDyslexic (Abbie Gonzalez). As licenças estão em \ui{Sobre} → \ui{Licenças} no
aplicativo.

\end{document}
